\section{Experimental Design}

Table~\ref{tab:configs} defines the configurations used in the experiments. Figure~\ref{fig:pipeline-architecture} is the superset workflow; each configuration activates a specific path through that workflow. The identifier names the workflow variant, while prefixes such as ``Controlled C1'' or ``External C1'' name the campaign in which the same C1 variant is executed. Result-table labels such as ``C1 scanner only'' are descriptive aliases for the same non-blocking scanner configuration.

The configuration set is an ablation design, not a taxonomy prescribed by a standard. The choices isolate build/test overhead, scanner observation, scanner automation, policy enforcement, remediation with rescan, and the combined scanner-remediation-policy path. This follows benchmarking guidance that security-tool evaluations should use explicit targets, comparable configurations, and outcome criteria \cite{pariziBenchmark2018}; prior empirical work showing that scanner warnings require workflow context \cite{johnsonStaticAnalysis2013,bellerStaticAnalysis2016}; secure-development guidance that emphasizes integrated CI/CD security activities \cite{nistSP800204D,nistSSDF,owaspSAMM}; and dependency-remediation studies showing that automated updates must be evaluated for acceptance and residual risk rather than treated as sufficient by construction \cite{decanNpm2018,mirhosseiniPullRequests2017}.

\begin{table*}[t]
\centering
\footnotesize
\caption{Pipeline configurations and active workflow stages.}
\label{tab:configs}
\begin{tabular}{@{}lp{0.43\linewidth}p{0.43\linewidth}@{}}
\toprule
ID & Active workflow path & Stages intentionally absent or interpretation \\
\midrule
C0 & Codebase copy and build/test mode only. & No security scanner, normalizer output, remediator, post-remediation rescan, or policy gate. \\
C1 & Codebase copy, initial scanners, normalizer, and run metrics. & No remediation, no post-remediation rescan, and no OPA gate. Residual findings equal initial findings; delivery is not blocked. \\
C2 & Codebase copy, automatic scanner mode, normalizer, and run metrics. & No remediation and no OPA gate; this isolates scanner automation behavior from policy enforcement. \\
C3 & Codebase copy, initial scanners, normalizer, raw-finding policy gate, and run metrics. & No remediation and no post-remediation rescan; OPA decides over raw findings. \\
C4 & Codebase copy, initial scanners, normalizer, bounded remediation, post-remediation rescan, and run metrics. & No OPA gate; this isolates whether remediation reduces residual findings before delivery policy is applied. \\
C5 & Codebase copy, initial scanners, normalizer, bounded remediation, post-remediation rescan, OPA gate, and run metrics. & Full SecFlowOps path; policy may deny delivery even when the experimental run succeeds. \\
B1 & NodeGoat package-manager remediation using \texttt{npm audit fix --package-lock-only --omit=dev}, followed by npm audit metrics. & No SecFlowOps secret remediation and no OPA gate; used only as a package-manager baseline. \\
\bottomrule
\end{tabular}
\end{table*}

\begin{table*}[t]
\centering
\footnotesize
\caption{Rationale for the selected configurations.}
\label{tab:config-rationale}
\begin{tabular}{@{}lp{0.34\linewidth}p{0.54\linewidth}@{}}
\toprule
ID & Evaluation role & Cited support \\
\midrule
C0 & Non-security execution baseline for build/test overhead. & Required to separate pipeline overhead from security-tool overhead under comparable benchmark conditions \cite{pariziBenchmark2018}. \\
C1 & Scanner-only observation baseline. & Static-analysis findings need workflow context and cannot by themselves establish release acceptability \cite{johnsonStaticAnalysis2013,bellerStaticAnalysis2016}; C1 measures that observation layer alone. \\
C2 & Automatic scanner mode without remediation or policy. & Secure CI/CD guidance motivates automated security activities \cite{nistSP800204D,nistSSDF,owaspSAMM}; C2 isolates scanner automation from remediation and gate decisions. \\
C3 & Policy-only decision over raw findings. & Policy gates are part of integrated DevSecOps control, but C3 tests whether gate behavior alone explains delivery outcomes \cite{nistSP800204D,opaDocs}. \\
C4 & Remediation and rescan without policy. & Dependency and remediation work should be evaluated by residual effects, not only by update creation \cite{decanNpm2018,mirhosseiniPullRequests2017}; C4 isolates remediation from gate enforcement. \\
C5 & Full scanner-remediation-rescan-policy path. & Integrated DevSecOps and secure-development guidance motivate combining continuous security checks, remediation evidence, and release controls \cite{nistSP800204D,nistSSDF,owaspSAMM}. \\
B1 & npm-audit-only package-manager baseline for NodeGoat. & npm advisories and automated dependency upgrades are a known special case requiring separate evaluation \cite{npmAudit,decanNpm2018,mirhosseiniPullRequests2017}. \\
\bottomrule
\end{tabular}
\end{table*}

Table~\ref{tab:campaigns} summarizes the campaigns. The controlled campaign evaluates six configurations: build only, non-blocking scanning, automatic scanning, policy only, agents only, and the combined SecFlowOps configuration. Five generated codebases inject SAST, SCA, secret, Docker, and Kubernetes findings with explicit ground-truth labels. Each controlled codebase is evaluated three times, yielding 90 runs.

The full protocol campaign executes the complete local design: three codebase families, eight scenarios, six configurations, and three repetitions, yielding 432 runs. The eight scenarios are clean, vulnerable dependency, fake secret, reflected XSS, SQL injection, Docker misconfiguration, Kubernetes misconfiguration, and multi-layer. The 432-run protocol records the DAST scenario as ground truth, while executable ZAP coverage is measured in a separate targeted DAST probe to avoid merging a focused runtime validation with the full local protocol.

The external campaign evaluates five configurations, non-blocking scanning, automatic scanning, policy only, agents only, and the combined SecFlowOps configuration, on five open-source codebases: Requests, Flask, Express, OWASP NodeGoat, and DVNA. Each external codebase is evaluated three times, yielding 75 runs. The external campaign uses an explicit build-skip mode to avoid confounding security-pipeline behavior with project-specific dependency installation and test setup.

Two focused campaigns address external validity gaps. A NodeGoat npm study repeats all five configurations three times on OWASP NodeGoat and enables \texttt{npm audit fix --package-lock-only --omit=dev} in the isolated remediation copy. A separate npm-audit-only baseline repeats the same package-manager remediation three times without SecFlowOps secret remediation or OPA gating. An injected external recall study copies Requests, Flask, and Express, injects six controlled findings, and evaluates three configurations over three repetitions, yielding 27 runs with complete injected ground truth.

\begin{table*}[t]
\centering
\footnotesize
\caption{Evaluation campaigns and main purpose.}
\label{tab:campaigns}
\begin{tabular}{@{}p{2.7cm}p{3.2cm}p{1.1cm}p{1.1cm}p{6.0cm}@{}}
\toprule
Campaign & Corpus & Runs & Reps & Purpose \\
\midrule
Controlled & five generated API variants & 90 & 3 & End-to-end behavior with injected ground truth and controlled size variation. \\
Full protocol & 24 generated scenario codebases & 432 & 3 & Complete local design: three families, eight scenarios, six configurations, and three repetitions. \\
External OSS & Requests, Flask, Express, NodeGoat, DVNA & 75 & 3 & Scanner, remediation, adjudication, policy, and performance behavior on real codebases. \\
NodeGoat npm & OWASP NodeGoat & 15+3 & 3 & Scanner-only, npm-audit-only, agents-only, and SecFlowOps behavior on a difficult npm dependency graph. \\
Injected external & Requests, Flask, Express copies with injected findings & 27 & 3 & Recall and residual target findings where complete injected ground truth exists in real-codebase contexts. \\
ZAP DAST probe & reflected-XSS generated codebase & 2 & 1 & Validate OWASP ZAP execution, JSON normalization, and DAST ground-truth matching. \\
\bottomrule
\end{tabular}
\end{table*}
